% chapters/06_fazit.tex

\chapter{Fazit und Ausblick}
\label{chap:fazit}

\section{Fazit}
\label{sec:fazit}

Diese Projektarbeit hat gezeigt, dass das Fine-Tuning des \ac{vla}-Modells
NVIDIA GR00T N1.6 auf dem humanoiden Roboter Unitree G1 mit DEX3-Hand für
Block-Stacking-Aufgaben prinzipiell möglich ist.

Die zentralen Erkenntnisse lassen sich wie folgt zusammenfassen:

\begin{itemize}
  \item Die entwickelte Docker- / Apptainer-Infrastruktur ermöglicht ein
    reproduzierbares Training lokal und auf dem \ac{hpc}-Cluster KISSKI.
  \item \ac{il} mit Behavior Cloning konvergiert auf dem gegebenen Datensatz
    innerhalb von~\TODO[Schritte bis Konvergenz]~Schritten auf eine stabile Verlustfunktion.
  \item Die Erfolgsrate von~\TODO[Erfolgsrate in \%]~\unit{\percent} in der Simulation liegt
    [über / unter / im Bereich von] vergleichbaren Arbeiten.
  \item Die Pipeline ist modular aufgebaut und lässt sich auf weitere
    Manipulationsaufgaben übertragen.
\end{itemize}

\section{Ausblick}
\label{sec:ausblick}

Auf Basis dieser Arbeit ergeben sich folgende Anschlussprojekte und
offene Fragen:

\paragraph{\ac{rl}-Verfeinerung}
Ein \ac{rl}-Fine-Tuning-Schritt auf Basis des \ac{il}-Checkpoints könnte die
Robustheit gegenüber Störungen und Variationen in der Szene erhöhen.

\paragraph{Sim-to-Real-Evaluation}
Eine systematische Evaluation auf dem physischen G1-Roboter steht noch aus
und wird als nächster Meilenstein des Projekts angestrebt.

\paragraph{Erweiterung des Datensatzes}
Die Kombination von menschlichen und synthetischen Demonstrationen aus der
NVIDIA-Isaac-Lab-Simulation könnte die Datenbasis erheblich vergrößern und
die Generalisierungsfähigkeit des Modells verbessern.

\paragraph{Multi-Embodiment-Training}
GR00T N1.6 unterstützt prinzipiell das gleichzeitige Training auf mehreren
Roboter-Plattformen; ein gemeinsames Modell für G1 und weitere humanoide
Systeme wäre ein langfristiges Ziel.
